\section{Related Work}

\subsection{Reinforcement Learning from Execution Feedback for Code Generation}

The use of execution signals to improve code language models is well-established. CodeRL \cite{Le2022CodeRL} pioneered the approach using binary pass/fail rewards in an actor-critic framework, demonstrating consistent improvements over SFT on APPS and HumanEval (+4.3\% pass@1). PPOCoder \cite{Majeed2023PPOCoder} applied PPO with binary execution reward achieving 5--8\% improvements on HumanEval and MBPP. RLTF \cite{Liu2023RLTF} introduced partial credit reward (fraction of tests passing) and reported approximately 7\% improvement on HumanEval/MBPP, with the claim that partial reward advantages are especially strong at harder problem instances.

These works share a common limitation: evaluation is confined to HumanEval and MBPP, covering only easy to medium difficulty. The critical question of whether RLEF advantages \textit{scale with benchmark difficulty} into hard competitive programming regimes remains unaddressed.

The most directly related prior study is RLEF \cite{Gehring2024RLEF}, which reports that partial reward outperforms binary reward more at harder problems. However, their study uses a proprietary Meta internal model, making independent verification impossible. Our work fills this reproducibility gap using publicly available components throughout.

\subsection{Reward Formulation in RLEF}

The choice between binary and fractional execution reward has received focused study. Two recent papers challenge the common intuition that fractional reward's denser signal should outperform binary, especially at hard difficulty. First, \cite{Anonymous2025PassRate} finds that 97\% of APPS problems produce identical solve/fail outcomes regardless of reward formulation at GRPO convergence. Second, VeRPO \cite{Anonymous2025VeRPO} identifies a \textit{cardinality bias} in na\"{i}ve fractional reward on variable-test-count datasets like APPS.

Our null result (p=0.552) is consistent with both of these papers on a different model/dataset combination. The DeepSeek-R1 work \cite{DeepSeekAI2025R1} further corroborates that binary execution feedback suffices at scale in mathematical reasoning.

\subsection{Hard Benchmark Evaluation for Code LLMs}

LiveCodeBench \cite{Jain2024LiveCodeBench} provides contamination-free evaluation stratified by difficulty from competitive programming contests. MBPP \cite{Austin2021MBPP} and HumanEval \cite{Chen2021HumanEval} are established but cover only easy/medium function-level problems. APPS \cite{Hendrycks2021APPS} is the most widely used algorithmic training dataset. Our coverage analysis reveals that 85.32\% of APPS competition problems have reference solutions --- refuting the ``dataset void'' framing used in prior mechanistic explanations.

\subsection{Our Position}

No prior work provides a controlled, reproducible, open-source comparison of RLEF vs.\ SFT across the full difficulty spectrum. Our work fills the gap between non-reproducible large-scale studies \cite{Gehring2024RLEF,DeepSeekAI2025R1} and easy/medium-only evaluations \cite{Le2022CodeRL,Majeed2023PPOCoder,Liu2023RLTF}.
